\begin{afterword}
\summaryhead

The author recalls reading the news of the September 11 attacks and thinking at once that the overreaction would be ten times worse than the attack. In the event, the post says, that was a vast understatement.

The reason given is a ``spiral of hate,'' the mirror of the happy death spiral. Everyone gains by stressing how terrible the attack was, and no one dares urge a proportionate response, so the reaction overshoots whatever the right level is. Attacking the Enemy makes you a patriot, and questioning any negative claim about the Enemy makes you a traitor. The example is the claim that the hijackers were cowards, which the author calls plainly untrue.

The post then says that the American response did the United States more harm than any terrorist group could, that ignoring the attack would have been better, and that politicians had no choice but to walk into al-Qaeda's trap. After a brief moment of candor, restraint became unspeakable within two days, and from then on fury and folly could only rise.

\responsehead

The record supports much of the post's description. On the day of the attacks the President spoke of ``cowardly acts.'' Three days later the House authorized force by 420 votes to 1, and its one dissenter was called a traitor and received death threats. By March 2008 the war in Iraq had killed more American soldiers than the attacks killed people. Bin Laden himself boasted in 2004 of having baited the administration.

The post goes beyond the record where it turns that description into absolutes. ``No one would dare to be the voice of restraint'': Barbara Lee did, in nearly those words, on 14 September. The claim that the hijackers were not cowards had been made in public within two weeks, by Bill Maher and Susan Sontag, as the notes on ``Belief as Attire'' record. The closing law is conditional: fury and folly ``can only rise with time'' once restraint becomes unspeakable. In the case the post describes, restraint had become speakable again by the week the post appeared, when most Americans polled by Gallup called the war in Iraq a mistake. The post describes how the spiral began and says nothing of how it slowed.

Its largest claims are the least argued. ``Ten times worse'' and ``a vast understatement'' come with no measure. That ignoring the attack ``would have been better than the real course of history'' is stated without a hedge, and considers only the costs of the actual response. The post names no war and no policy, so the reader must supply what ``shooting off its own foot'' refers to.

In short: a mechanism that the record of September 2001 largely supports, stated in an absolute (no one dared) that the record does not support, with its largest claims given no measure and no account of how the spiral slowed.
\end{afterword}
